\section{Discussion}

The construction produces two canonical global structures on the
Petersen line-graph core:

\begin{itemize}
\item the \emph{Thalean cocycle}, a nontrivial
$\mathbb{Z}_2$-valued invariant arising from transport parity in the
double cover $G_{30}\to G_{15}$;

\item the \emph{Thalean sector system}, whose incidence operator $M$
satisfies the polynomial identity
\[
M M^{\mathsf T} = A^3 + 2A^2 + 2I,
\]
giving a rigid algebraic and geometric structure.
\end{itemize}

These structures are not independent: both arise from the same lifted
transport on the dodecahedral flag space. The cocycle records the
holonomy of transport, while the sector system encodes its cumulative
incidence geometry.

Together, they show that local transport rules on a combinatorial
polyhedron can generate global algebraic invariants and highly
constrained spherical configurations on associated quotient graphs.

This suggests a general mechanism: transport systems on flag spaces of
polyhedra may systematically induce cohomological data and polynomial
representations of adjacency algebras on their core quotients.

\subsection*{Acknowledgements}

The author acknowledges the use of AI-assisted tools in the development,
structuring, and refinement of this manuscript.
